El Banco Central de Guatemala requiere monitorear la evolución de los principales indicadores monetarios y financieros para comprender el comportamiento de la liquidez, el crédito al sector privado, las remesas familiares y la actividad económica del país. El objetivo es generar información que apoye la toma de decisiones relacionadas con la estabilidad financiera y el seguimiento del crecimiento económico.